\subsection*{b.}
\addcontentsline{toc}{subsection}{Inciso b}

\textbf{b. Supón que una empresa en crecimiento necesita implementar un sistema de administración de bases de datos (SMBD) para centralizar su información operativa (ventas, clientes, inventario y logística). Elige dos posibles soluciones (p.e., PostgreSQL, Oracle, MySQL, SQL Server, etc.) y compara sus ventajas y desventajas en términos de: rendimiento ante grandes volúmenes de datos; escalabilidad a futuro; costos de licenciamiento y mantenimiento; soporte técnico y comunidad. Con base en tu análisis, argumenta cuál solución recomendarías y por qué.}

Para realzar este análisis, se realizará una comparativa entre los sistemas manejadores de bases de datos (SMBD): PostgreSQL y Oracle. Particularmente, como Oracle tiene una versión gratuita y una versión de pago, que no es fija y depende de las necesidades del sistema, para algunos criterios se tomaran en cuenta ambos y se compararán con PostgreSQL. El análisis se realizará en base a los siguientes 4 criterios.

\begin{itemize}
\item \textbf{Rendimiento del sistema ante grandes volúmenes de datos:} Tanto Oracle en su versión de pago y PostgreSQL están diseñadas para operar con grandes volúmenes de datos. Cabe mencionar que la versión gratuita de Oracle está limitada en cuanto a los recursos disponibles, teniendo un límite de almacenamiento de 12 GB.
\item \textbf{Escalabilidad: } Del mismo modo, Oracle en su versión de pago y PostgreSQL, son buenas opciones, ya que, por un lado, las licencias de pago de Oracle, permiten que el sistema sea tan amplio como la necesidad lo requiera, mientras que a nivel de Software, PostgreSQL, no ofrece ninguna limitante en cuanto a memoria RAM o número de procesadores que pueden ser destinados al mantenimiento del sistema, por lo que ambos sistemas son escalables. Cabe mencionar que la versión gratuita de Oracle impone limitaciones considerables a nivel de recursos, permitiendo usar un máximo de 2 procesadores y 2 GB de RAM.
\item \textbf{Costos de licenciamiento y mantenimiento: } Ambos requieren un costo de infraestructura para el almacenamiento de los datos; sin embargo, en Oracle se necesita pagar una licencia para sostener su funcionamiento con el aumento de los datos y una licencia de mantenimiento para mantener el software actualizado, mientras que PostgreSQL es un software completamente gratuito.\
\item \textbf{Soporte técnico y comunidad.} Con la licencia de pago de Oracle, se ofrece un soporte técnico especializado en resolver cualquier problema con el manejo del software, mientras que PostgreSQL, al ser de código abierto, no ofrece un soporte técnico directo, pero cuenta con una amplia comunidad para solucionar cualquier problema que pueda surgir; el soporte y mantenimiento del sistema queda en manos de la empresa.
\end{itemize}
Del análisis anterior, consideramos que Oracle, con su licencia de pago, presenta ventajas considerables al momento de realizar un mantenimiento al sistema y para solucionar problemas técnicos; sin embargo, por el costo elevado de sus licencias, consideramos que la mejor opción para una empresa en crecimiento es PostgreSQL, ya que no requiere un costo extra por manejo de software, ofrece una gran escalabilidad y tiene una gran capacidad para sostener los datos originados por la empresa con su crecimiento.

